\documentclass[10pt]{article}
\usepackage[margin=20mm]{geometry}
\usepackage{amsmath,amssymb,mathtools,fontspec,unicode-math,xurl,hyperref}
\IfFontExistsTF{DejaVu Serif}{\setmainfont{DejaVu Serif}}{\setmainfont{Latin Modern Roman}}
\setmathfont{latinmodern-math.otf}
\setlength{\parindent}{0pt}\setlength{\parskip}{6pt}\sloppy
\title{Initial scientific review: M25-0033}
\author{}
\date{25 October 2026}
\begin{document}
\maketitle
\section*{Identity, version, and decision}
Paper: Tongmu He, Perfectoidness via Sen theory and applications to Shimura varieties, DOI 10.1090/jams/1060. The reviewed object is the 59-page author manuscript arXiv:2407.14488v2 (arXiv submission 23 July 2025; internal manuscript date 24 July), with supplied source00000TOKEN. Its PDF SHA-256 is\par
> 99058dd978553cdd5d98be85453619d8e2e1
> c0861aad77d78028e0e4242a626f.\par
Despite the source filename, no comparison with the published journal version has been performed. Every assessment below concerns this exact author manuscript.\par
Judgment: No Issue Found within the inspected scope. No demonstrated invalid derivation or counterexample was established in the checks below. This is a bounded initial review, not certification of all 59 pages or of the external foundations. In particular, the construction of the Faltings extensions, several ramification estimates, and the full comparison with geometric Sen operators remain outside the independent verification completed here.\par
\section*{Actual coverage and dependency map}
The text of PDF pages 1–6, 16–18, 39–44, 46–51, and 56–57 was read. Images of pages 1, 3–6, 16–18, 40–41, 43–44, 46–51, and 56–57 were inspected for the statements and formulas used below. Critical occurrences in the supplied TeX were checked, including lines 284, 362, 1159–1161, 2250–2252, 2640, 2976, and 3037. Extraction of the full PDF and hashing the supplied sources do not constitute a complete semantic review.\par
The inspected core chain is \[\begin{aligned}
 &\text{geometric Sen injectivity}
 \Longrightarrow \text{full Sen injectivity}\\
 &\quad\Longrightarrow \text{maximal invariant dimension}
 \Longrightarrow \delta=0\\
 &\quad\Longrightarrow \text{pointwise pre-perfectoidness}.
 \end{aligned}\] Boundary root divisibility is then needed for the open variety; the almost comparison and a dimension bound imply integral torsion-cohomology vanishing. The Shimura application additionally imports the locally split injectivity theorem quoted as Theorem 11.2. The checks below distinguish deductions verified here from results used conditionally.\par
\section*{Definitions, quantifiers, and boundaries}
The relevant hypothesis in Theorem 8.9 is injectivity of the geometric Sen action at the specified geometric valuative point, not merely that the map is nonzero. For geometric dimension greater than one, a nonzero linear map can have a kernel. Corollary 9.13 requires its valuative-point hypothesis for every point in the stated range, and separately for each generic component with the stipulated Galois extension. One favorable point cannot establish the global conclusion.\par
The distinction between a pre-perfectoid field and its completed perfectoid field is retained throughout this report. The inspected statements establish the stated valuation-theoretic conclusions; they are not, by themselves, a proof that a chosen global infinite-level space is a perfectoid space.\par
An apparent ground-field concern was rejected after checking the original image of page 4 and TeX line 362: Theorem 1.19 is formulated over \(\overline K\). The overline is lost in some plain-text extraction. A proposed arithmetic example over \(K\) instead of \(\overline K\) therefore does not satisfy the statement’s hypotheses and is not a counterexample. This is recorded as counterevidence, not as a paper defect.\par
The dimension-zero case is consistent with the inspected linear algebra: the geometric part has dimension zero while the arithmetic contribution still has to be nonzero. When the normal-crossings divisor is empty, the requirement to divide its local equations is vacuous. No vanishing at the boundary degree \(q=\dim Y\) is asserted.\par
\section*{Supporting checks}
\subsection*{The Lie bracket in Section 8.2 (page 41)}
Let \(E\) be the relevant vector space and let \(\lambda=\iota^*\) be the linear functional used in the construction. The proposed bracket is \[[x,y]=\lambda(x)y-\lambda(y)x .\] It is alternating and bilinear, and \(\lambda([x,y])=0\). Its iterated bracket is \[[[x,y],z]=-\lambda(z)\lambda(x)y+\lambda(z)\lambda(y)x .\] The cyclic sum cancels term by term, proving Jacobi. Thus \(\ker\lambda\) is an abelian ideal. If \(\lambda(e)=1\), then \([e,v]=v\) for \(v\in\ker\lambda\), giving the advertised semidirect-product description. This verifies the algebraic step conditional on the previously constructed extension and functional; it does not reconstruct that extension.\par
\subsection*{Invariant dimension and the connecting homomorphism}
In Lemma 4.22 (page 17) and Lemma 7.18 (page 40), the invariant sequence has the relevant shape \[0\longrightarrow \widehat F(1)\longrightarrow E^G
 \xrightarrow{j} V\xrightarrow{\delta} H^1(G,\widehat{\overline F}(1)),
 \qquad \dim_{\widehat F}V=d .\] The hypotheses fixing the necessary roots of unity and the identification of invariant scalars are essential. Under those hypotheses the left term has dimension one. Exactness gives \[\dim E^G=1+\dim\ker\delta\leq 1+d.\] Consequently maximal dimension is equivalent to surjectivity of \(j\), and to \(\delta=0\). This argument uses the connecting homomorphism and left exactness; it does not assume that taking invariants is right exact. The same computation remains valid for \(d=0\).\par
\subsection*{The trace contradiction in Theorem 4.23 (pages 17–18)}
The proof uses torsion-freeness of \(M=\Omega^1_{\mathcal O_F/\mathcal O_{\overline K}}\) and the previous normalized-trace results. Conditional on torsion-freeness, \(M\) embeds in its fraction-field differential space. If a transcendence-basis differential \(dt_i\) is infinitely \(p\)-divisible in \(M\), then for any \(a\in F\) the height-one valuation permits an integer \(n\) with \(p^n a\in\mathcal O_F\), and \[a\,dt_i=(p^na)\,(dt_i/p^n)\in M.\] If all basis differentials have this property, \(M\) equals the indicated \(F\)-linear differential space.\par
In the contrary case the Kummer cocycle \(\xi_t\) represents the connecting class. Vanishing of that class gives a coboundary \(x\). The kernel of the Kummer character and the invariant-field identification place \(x\) in the completed Kummer extension. The normalized trace \(T\) is an equivariant retraction to \(\widehat F\). Hence the nonzero value \((\tau-1)x=\xi_t(\tau)\) would lie both in \(\widehat F\) and in \(\ker T\), impossible since their intersection is zero. This checks the final contradiction using the precise retraction property. It is conditional on the established trace estimates and on the cited differential/pre-perfectoidness results. Remark 4.24 explicitly says that stronger estimates used for the homeomorphism in 4.19(3) are not required here; they were not treated as dependencies of this particular deduction.\par
\subsection*{Theorem 8.9: geometric to full injectivity (pages 43–44)}
The diagram in 8.8 has geometric and arithmetic exact rows. The cyclotomic hypothesis ensures that the arithmetic map is injective on its one-dimensional source. If \(x\) is in the kernel of the full action, its arithmetic image vanishes; injectivity of the arithmetic map forces \(x\) into the geometric subspace, where geometric injectivity forces \(x=0\). This is a direct diagram chase.\par
Dualizing an injective map of finite-dimensional vector spaces gives a surjection to the Faltings extension. Under the stated Galois and cyclotomic hypotheses the appropriate source basis is invariant. Its images are invariant and span the target. This proves the required invariant-dimension conclusion without incorrectly asserting that an arbitrary surjection remains surjective after applying invariants. The dimension count also requires the full target Lie algebra to have dimension at least \(d+1\) in this situation. The implication from the resulting zero connecting map to pre-perfectoidness and to boundary \(p\)-power roots uses 7.19 and 7.20. The last Kummer contradiction in 7.20 was followed; the entire earlier construction was not independently reproved.\par
\subsection*{Corollary 9.9: from almost vanishing to actual vanishing}
Here the ground field in 9.6 is explicitly an algebraically closed valuation field of height one (page 46). The cohomological comparison 9.7 and the pre-perfectoid calculation 9.8 are used conditionally. The subsequent dimension and coefficient arguments were checked.\par
For a proper, flat, finitely presented model, fiber dimension is locally constant; the valuation base is connected. The invoked Stacks Project lemmas 0D4J and 0E7D have the appropriate flatness, finite-presentation, properness, and quasi-coherent hypotheses. For \(q\) exceeding the fiber dimension the higher direct image vanishes, and the affine base has no further quasi-coherent cohomology. The passage along the directed system uses the cited continuity results.\par
There is no general principle that an almost zero module must be zero. The particular module here is \(V\otimes_{\mathbf F_p}\mathcal O_L/p\), where \(V\) is an \(\mathbf F_p\)-vector space. If \(V\ne0\), choose a basis vector; its span gives a direct summand isomorphic to \(\mathcal O_L/p\). This module is not almost zero: since \(L\) is algebraically closed, choose \(\pi\) with \(0<v(\pi)<v(p)\). Then \(\pi\in\mathfrak m_L\) but \(\pi\cdot1\ne0\) modulo \(p\). Thus the asserted almost vanishing forces \(V=0\).\par
The coefficient induction also preserves the degree bound. The short exact sequence with successive quotients \(\mathbf Z/p\) gives a long exact sequence; if both adjacent \(H^q\) groups vanish for \(\mathbf Z/p^{n-1}\) and \(\mathbf Z/p\), then \(H^q\) with \(\mathbf Z/p^n\) vanishes. The completed-cohomology definition in the introduction is an ordinary inverse limit of these already zero direct limits. No interchange of derived inverse limit and cohomology, or unproved \(\varprojlim^1\) vanishing, is needed for this final step.\par
\subsection*{The Shimura application (pages 56–57)}
Theorem 11.2 is an external input asserting locally split injectivity. Its proof was not independently verified. Conditional on it and the comparison in 10.4–10.5, the application gives the stated vanishing over \(\overline{\mathbf Q}_p\).\par
The transfer to \(\mathbf C\) uses invariance of torsion etale cohomology under extensions of separably closed fields (Stacks 0F0B). At finite level the varieties are quasi-projective and hence quasi-compact and quasi-separated; the torsion order is invertible in characteristic zero. One can compare both fields through an algebraic closure of the reflex field with compatible embeddings, and then pass to the specified filtered colimit. The fields need not be canonically embedded into one another. These are the hypotheses behind the short sentence in the proof; no violation of the cited base-change theorem was found.\par
\section*{Findings and counterevidence}
There are no confirmed numbered defects in this inspected portion. The following rejected or bounded concerns are retained to prevent overstatement:\par
1. A counterexample over the unclosed arithmetic field does not meet Theorem 1.19’s printed \(\overline K\) hypothesis.\par
2. Nonzero geometric Sen action is weaker than the precise injectivity hypothesis. The report evaluates the latter.\par
3. Almost vanishing is upgraded to vanishing only for the displayed scalar extension of an \(\mathbf F_p\)-vector space; the explicit direct-summand argument justifies that upgrade.\par
4. The final Shimura conclusion remains dependent on external Sen and comparison theorems. Checking the formal application does not verify their proofs.\par
None of these items supplies a demonstrated Major or Minor defect. No negative judgment is inferred from the amount of unreviewed material.\par
\section*{Version search and remaining verification}
The official arXiv record was checked for this exact identifier and version. A bounded title/erratum search did not identify a correction relevant to the inspected steps. This is not a claim that no erratum exists. No journal-final PDF was compared.\par
Remaining work includes a complete semantic review of the uninspected pages; the construction and functoriality of Faltings extensions; the preceding ramification estimates and adequate-algebra arguments; the full boundary and cohomological comparison proofs; the exact cited results of the author’s earlier papers and Rodriguez Camargo; and journal-version comparison. The report does not recommend unconditional scientific acceptance on this limited evidence.\par
\section*{Supporting checks}
Judgment | No Issue Found within the inspected scope.\par
Repair | No demonstrated defect requiring a repair was established.\par
Verification | Explicit algebra, diagram, dimension, almost-module, coefficient, and base-change checks above; selected original pages and TeX inspected.\par
Adoption | Bounded initial report only; full-paper acceptance remains unverified.\par
Remaining | Uninspected proofs, external foundations, and journal comparison.\par
\section*{Supporting checks}
Only this task’s staging directory was written. The input integrity pass verified the manifested sizes and hashes; this is not a proof audit. All scientific reasoning and source-page inspection in this report were performed by the main assistant. Author programs and author TeX were not executed.\par
These statements are kept separate; no runtime-equivalence assertion or personally observed browser-tab claim is made, and no independent backend mismatch was established.\par
The original environment probe preceded the authorized portable Tectonic build. The follow-up build and its page inspections are preserved as historical evidence. Original-review actions retain their original attribution. Focused AI acceptance checks the declared scope without extending it to the uninspected proofs. Scientific adjudication remains pending and no human review is claimed.\par
9 T. He, Perfectoidness via Sen theory and applications to Shimura varieties, arXiv:2407.14488v2, supplied 59-page author manuscript and TeX. https://arxiv.org/abs/2407.14488v2. The Stacks Project, Lemma 37.30.6, Tag 0D4J. https://stacks.math.columbia.edu/tag/0D4J. The Stacks Project, Lemma 32.19.2, Tag 0E7D. https://stacks.math.columbia.edu/tag/0E7D. The Stacks Project, Lemma 59.90.2, Tag 0F0B. https://stacks.math.columbia.edu/tag/0F0B.\par
\end{document}
